\documentclass[11pt,a4paper]{article}
\usepackage[margin=24mm]{geometry}
\usepackage{amsmath,amssymb,amsthm}
\setlength{\emergencystretch}{2em}
\newtheorem{theorem}{Theorem}
\newcommand{\gen}[1]{\langle#1\rangle}
\title{A third fixed point in a singleton prefix}
\author{Human Compression project: L044 continuation}
\date{}
\begin{document}\maketitle
Retain the complete critical singleton hypotheses and the acquired
small-support budget. Thus $n=2r\ge14$, the actual natural
$G_0$ is $\operatorname{Alt}(\Omega\setminus\{a\})$, full ambient IP
holds, and $G$ is a transitive even string action. For $4\le g<r$ put
\[
 P=G_{<g},\quad Y=\operatorname{supp}(P),\quad C=a^P,
 \quad H=G_{\{2,\ldots,g-1\}},\quad b=a^{\rho_0}.
\]
Assume $P$ is intransitive on $Y$, $|Y|\le2g+1$, and every
colour $2,\ldots,g-1$ has two ambient-eligible edges.
\begin{theorem}
The group $H$ cannot have precisely three fixed points in $Y$
and precisely two nontrivial orbits.
\end{theorem}
\begin{proof}
Suppose this decomposition holds. The acquired budget gives
$|Y|=2g+1$, fixed set $\{a,b,z\}$ and two nontrivial orbits
$U,V$. On the fixed set $\rho_0=(a\ b)$ and fixes $z$.
Since $z\in Y$ and both $\rho_0$ and $H$ fix it,
$\rho_1$ moves $z$.

Strict growth gives $|C|\ge g+1\ge5$, so $C$ contains an
$H$-orbit. It cannot contain both: the only possible point left
outside would be $z$, which is moved by $P$. Thus $P$ has exactly
two nontrivial orbits in $Y$, namely
\[
 C=\{a,b\}\cup V\cup Z_C,\qquad E=U\cup Z_E,\qquad
 \{Z_C,Z_E\}=\{\{z\},\varnothing\}.                 \tag{1}
\]
Both are preserved by $\rho_1$. The involution $\rho_0$
centralizes $H$. It cannot exchange $U,V$, since then the
$P$-orbit $C$ would contain both. On each preserved orbit it is
either free or the identity. It cannot be the identity on both,
since it would be the odd permutation $(a\ b)$.

If it is free on both, its support is $Y\setminus\{z\}$.
Since $\rho_1$ moves $z$, the group
$P_1=\gen{\rho_0,\rho_1}$ has support exactly $Y$.
Every mark of index at least three centralizes $P_1$ and preserves
that support. The mark $\rho_2\in P$ also preserves $Y$.
Thus $Y$ is a proper nonempty $G$-invariant set, impossible.
Consequently $\rho_0$ is free on exactly one of $U,V$.
Its evenness makes the size of that orbit twice an odd integer;
the acquired exclusion of an internally swapped two-point
$H$-orbit makes its size at least six.

We shall repeatedly use the actual first-prefix stabilizer.
Write $C_1=a^{P_1}$ and $m=|C_1|$. Strict growth gives $m\ge3$.
Actual natural $G_0$ and full IP imply
\[
 (P_1)_a=P_1\cap G_0=\gen{\rho_1},\qquad |P_1|=2m.       \tag{2}
\]
Hence the product of the two distinct involutions has order $m$.
Every mark of index at least three centralizes $P_1$, fixes $a$,
and therefore fixes $C_1$ pointwise.

Put $K=G_{\{3,\ldots,g-1\}}$. This is a nonempty interval of
nonidentity marks. Its support lies in $Y$, and its colours have
ambient-eligible pairs. If $f_K$ and $c_K$ count its fixed points
in $Y$ and its nontrivial orbits, the acquired interval theorem gives
\[
 2(g-3)\le |Y|-f_K-c_K,
 \qquad f_K+c_K\le7,\qquad c_K\ge1.                   \tag{3}
\]

\emph{First suppose $\rho_0$ is free on $U$ and fixes $V$.}
The point $x=b^{\rho_1}$ is different from $b$, lies in $C$
and is fixed by $\rho_0$. Thus $C_1=\{a,b,x\}$, with
marks $(a\ b)$ and $(b\ x)$. Equation (2) gives actual
$P_1\cong S_3$. On $C\setminus C_1$, the mark $\rho_0$
is trivial, so the relation $(\rho_0\rho_1)^3=1$ also makes
$\rho_1$ trivial there.

If $z\in C$, then $E=U$ is free under $\rho_0$.
Its $P_1$-orbits have sizes two or six, since their sizes divide
six and cannot be odd. They are precisely all the orbits of those
sizes in the ambient set. Every mark of index at least three
preserves their union $E$, while $\rho_0,\rho_1,\rho_2$ do
so by (1) and the $H$-orbit property. This contradicts transitivity.

Suppose instead that $z\in E$. Its $P_1$-orbit $D$ has size
three. Indeed it is nontrivial because $\rho_1$ moves $z$;
the possible orbit sizes are two, three and six. On a two-point
orbit $\rho_0$ has either zero or two fixed points, and a
six-point transitive action of a group of order six is regular.
Neither can contain the unique $\rho_0$-fixed point $z$ in $E$.
Every other orbit in $E$ has size two or six. Thus $C_1,D$ are
the only two three-point $P_1$-orbits globally. Each mark of index
at least three fixes $C_1$ pointwise, hence preserves $D$.
It also fixes $D$ pointwise: it centralizes $P_1$ and must fix
the unique $\rho_0$-fixed point $z$ in $D$, after which
transitivity propagates fixation.

Consequently $f_K\ge6$. Equation (3) forces $f_K=6,c_K=1$.
The sole nontrivial $K$-orbit is $Y\setminus(C_1\cup D)$,
of odd size $2g-5$. Since $\rho_0$ centralizes $K$, it
preserves this unique nontrivial orbit and fixes it pointwise:
a centralizing involution on an odd transitive orbit has a fixed
point, then fixes the whole orbit. But $D$ contains only two
points of $U$, and $|U|\ge6$. The remaining points of $U$
lie in that odd orbit and are moved by $\rho_0$, a contradiction.

\emph{Now suppose $\rho_0$ is free on $V$ and fixes $U$.}
If $z\in E$, then $\operatorname{supp}(\rho_0)=C$.
The mark $\rho_1$ preserves $C$, and every later mark centralizes
$\rho_0$, so $C$ is globally invariant, impossible.
We are left with $z\in C$, $E=U$ and
\[
 \operatorname{supp}(\rho_0)=C\setminus\{z\},
 \qquad \rho_0|_E=1.                                  \tag{4}
\]
Both marks preserve $C,E$, and $\rho_1$ moves $z$.

If $m$ is odd, the relation $(\rho_0\rho_1)^m=1$ and
$\rho_0|_E=1$ imply $\rho_1|_E=1$. By (4) and the movement
of $z$, the support of $P_1$ is then exactly $C$.
Every mark of index at least three preserves it by centralization;
$\rho_2\in H$ preserves it by (1). Again $G$ is intransitive.

If $m$ is even, consider $D=z^{P_1}\subseteq C$.
The involution $\rho_0$ has exactly one fixed point on $D$,
so $|D|$ is odd. As $\rho_1$ moves $z$, $|D|\ge3$.
Thus $D\ne C_1$. Every other $P_1$-orbit in $C$ has even
size, because $\rho_0$ is free on it; every orbit in $E$ has
size one or two, because $\rho_0$ is trivial there. Points
outside $Y$ are fixed. Therefore $D$ is the unique nontrivial
odd-sized $P_1$-orbit in the whole ambient set.
Each mark of index at least three preserves $D$, fixes its
unique $\rho_0$-fixed point $z$, and hence fixes $D$ pointwise.
It also fixes $C_1$ pointwise. Since even $m\ge3$ means $m\ge4$,
we obtain $f_K\ge |C_1|+|D|\ge7$, contrary to (3).
All cases are excluded.
\end{proof}

\paragraph{Attribution and remaining scope.}
The critical singleton setting, strict prefix growth, ambient interval
bound and internal two-point exclusion are acquired. The complete
support split, the actual first-prefix orbit argument, and the use of
six or seven fixed points to settle the third fixed point are this
continuation. The source context is Cameron--Fernandes--Leemans--Mixer,
\emph{Highest rank of a polytope for $A_n$}, Propositions~5.28--5.32.
Together with the separately accepted $f=2,c=2$ exclusion, this leaves
only $f=2,c=3$ in the stated small-support budget. No projected IP,
general faithful orbit action, or finite census is assumed. Complete
singleton routing and the other principal L044 obligations remain open
within this recovery; no whole-original or formal credit is added.
\end{document}
